\documentclass[a4paper,10pt,titlepage,oneside,openany]{book}

% ===== Pacchetti =====
\usepackage[skip=6pt]{parskip}
\usepackage[english]{babel}
\usepackage[utf8]{inputenc}
\usepackage[T1]{fontenc}
\usepackage{graphicx}
\usepackage{amsmath,amssymb}
\usepackage{amsthm}
\usepackage{enumitem}
\usepackage{float}
\usepackage{changepage} % per adjustwidth in pseudocode
\usepackage[top=1in, bottom=1in, left=1in, right=1in]{geometry}
\usepackage{fancyhdr}
\usepackage{tcolorbox}
\usepackage{array}
\usepackage{booktabs}
\usepackage{tabularx}
\usepackage{pdfpages}
\usepackage[colorlinks=true, linkcolor=black, citecolor=black, urlcolor=blue]{hyperref}
\usepackage{etoolbox}   % patch capitoli
\usepackage{ragged2e}   % per \justifying

% Pacchetti aggiuntivi necessari al corpo del paper
\usepackage{mathtools}
\usepackage{caption}
\usepackage{subcaption}
\usepackage[ruled,vlined]{algorithm2e}
\usepackage{listings}
\usepackage{xcolor}

% Compatibilita' con citazioni in stile \textcite senza biblatex
\newcommand{\textcite}[1]{\cite{#1}}

% Giustificazione e no-indent di default su tutto il documento
\AtBeginDocument{\justifying}
\setlength{\parindent}{0pt}
\setlength{\parskip}{6pt}

% -------------------------------
% Da cambiare quando si copia il template in \graphicspath{{img/}}
% -------------------------------
\graphicspath{{img/}}

% ===== Header / Footer =====
\pagestyle{fancy}
\fancyhf{}
\fancyhead[L]{\textbf{Starting Finance Club PoliTo}}
\fancyhead[R]{\includegraphics[width=1cm]{img/Logo_SFClub_Polito.png}}
\fancyfoot[C]{\thepage}
\renewcommand{\headrulewidth}{0.4pt}
\renewcommand{\footrulewidth}{0.4pt}
\setlength{\headheight}{25pt}
\setlength{\headsep}{24pt} % spazio sotto l'header

\fancypagestyle{plain}{%
  \fancyhf{}
  \fancyhead[L]{\textbf{Starting Finance Club PoliTo}}
  \fancyhead[R]{\includegraphics[width=1cm]{img/Logo_SFClub_Polito.png}}
  \fancyfoot[C]{\thepage}
  \renewcommand{\headrulewidth}{0.4pt}
  \renewcommand{\footrulewidth}{0.4pt}
}

% ===== Ambienti theorem ecc. =====
\theoremstyle{plain}\newtheorem{theorem}{Theorem}[chapter]
\theoremstyle{definition}\newtheorem{definition}{Definition}[chapter]
\newtheorem{notation}{Notation}[chapter]
\theoremstyle{remark}\newtheorem{remark}{Remark}[chapter]

\newenvironment{pseudocode}{\par\small\ttfamily\begin{adjustwidth}{1em}{1em}}{\end{adjustwidth}\par}

\newenvironment{tightabstract}[1][0.88]{%
  \begin{center}
  \begin{tcolorbox}[title={Abstract}, width=#1\textwidth, left=5mm, right=5mm,
    colback=white, boxrule=0.4pt]
  \small\justifying
}{%
  \end{tcolorbox}
  \end{center}
}

% Summary come chapter* + TOC + header, con noindent + giustificazione
\newenvironment{summary}{%
  \clearpage
  \chapter*{Summary}%
  \addcontentsline{toc}{chapter}{Summary}%
  \markboth{Summary}{Summary}%
  \noindent\justifying
}{}

\newenvironment{references}{\begin{thebibliography}{99}}{\end{thebibliography}}

\newcommand{\miniabstract}[1]{%
\vspace{-2pt}
\begin{tcolorbox}[colback=gray!6,colframe=black!15,boxrule=0.3pt,arc=3pt,left=6pt,right=6pt,top=4pt,bottom=4pt]
\textit{Mini-abstract.} #1
\end{tcolorbox}
\vspace{-4pt}
}

% ===== Fix capitoli: -18pt e dimensioni numeriche =====
\makeatletter
\patchcmd{\@makechapterhead}{\vspace*{50\p@}}{\vspace*{-18\p@}}{}{}
\patchcmd{\@makeschapterhead}{\vspace*{50\p@}}{\vspace*{-18\p@}}{}{}
\patchcmd{\@makechapterhead}{\huge}{\fontsize{17}{20}\selectfont\bfseries}{}{}
\patchcmd{\@makechapterhead}{\Huge}{\fontsize{20}{24}\selectfont\bfseries}{}{}
\patchcmd{\@makeschapterhead}{\Huge}{\fontsize{20}{24}\selectfont\bfseries}{}{}
\patchcmd{\@makechapterhead}{\vskip 20\p@}{\vskip 6\p@}{}{}
\patchcmd{\@makechapterhead}{\vskip 40\p@}{\vskip 10\p@}{}{}
\patchcmd{\@makeschapterhead}{\vskip 40\p@}{\vskip 10\p@}{}{}
\makeatother

% ===== Listings =====
\lstdefinestyle{python}{
  language=Python,
  basicstyle=\ttfamily\small,
  keywordstyle=\bfseries,
  commentstyle=\itshape,
  stringstyle=\ttfamily,
  showstringspaces=false,
  breaklines=true,
  frame=single,
  numbers=left,
  numberstyle=\tiny,
  xleftmargin=2em,
  framexleftmargin=1.5em
}
\lstset{style=python}

% ===============================
\begin{document}

% --- Copertina PDF come prima pagina, senza header/footer ---
\includepdf[pages=1,fitpaper=true,pagecommand={\thispagestyle{empty}}]{Copertina_research.pdf}

% --- TITLE/CONTACT/ABSTRACT ---
\begin{titlepage}
\thispagestyle{empty}
\vspace*{-1cm}
\begin{center}
\includegraphics[width=3.70cm]{img/Logo_SFClub_Polito.png}\\[1cm]
{\LARGE\bfseries High-Frequency Market Making under Avellaneda--Stoikov:\ Classical Diffusions and Multi-Asset Inventory Control\par}
\vspace{0.2cm}
{\large Single-asset benchmark models, covariance-aware extension, and event-driven simulation\par}

\vspace{0.8cm}

% --- Authors + reviewer + date (same size) ---
{\normalsize
\textbf{Authors:}\\[0.3cm]
Salvatore Gabriele Messina\\
Niccol\`o Soriano\\[0.3cm]
\textit{Date:} May 15, 2025\par
}

\vspace*{0.2cm}

% --- Only this line in Large (does not affect the rest) ---
{\Large\textbf{Starting Finance Club PoliTo}\par}

\vspace{0.4cm}
\normalsize

\textit{Contact:}\\
\small
LinkedIn: \url{https://www.linkedin.com/company/startingfinance-club-polito}\\
GitHub: \url{https://github.com/StartingFinanceClubPoliTo/Research}\\
Instagram: \url{https://www.instagram.com/sfclubpolito}\\
Website: \url{https://sfclubpolito.it/}

\vspace{1cm}
\begin{tightabstract}
This paper studies high-frequency market making within the Avellaneda--Stoikov framework under classical single-asset price dynamics and a covariance-aware multi-asset extension. We consider two finite-horizon benchmark cases: arithmetic Brownian motion with finite horizon and geometric Brownian motion with finite horizon. Across both cases, we compare a naive symmetric quoting policy with an inventory-based quoting policy derived from the Avellaneda--Stoikov control principle. The analysis is simulation-first and does not rely on market data. Instead, all experiments are conducted through a unified event-driven Monte Carlo engine that keeps dealer accounting and order-arrival mechanics fixed while varying the mid-price dynamics, and then extends the same logic to vector-valued prices and inventories. This design isolates the impact of inventory control, multiplicative price risk, and covariance-aware portfolio risk on dealer performance. The paper shows how inventory-aware quoting improves risk-adjusted performance by reducing inventory excursions and stabilizing marked-to-market wealth relative to naive quoting. The multivariate extension replaces the scalar inventory term $\sigma^2 q_t$ with the portfolio pressure $\Sigma q_t$, clarifying how correlated inventories modify quote centers and marginal bid/ask reservation prices. Overall, the study provides a coherent benchmark for understanding the role of classical diffusion modeling and covariance-aware inventory control in high-frequency market making.
\end{tightabstract}
\end{center}
\end{titlepage}

% --- Front matter: TOC in roman numerals starting from i ---
\frontmatter
\clearpage
\pagenumbering{roman}
\setcounter{page}{1} % "i"
\tableofcontents
\thispagestyle{plain}

% --- Main matter: body in arabic ---
\mainmatter
\clearpage
\setcounter{page}{1}

% ===============================
% Body
% ===============================

\chapter*{Introduction and Motivation}
\addcontentsline{toc}{chapter}{Introduction and Motivation}

\section*{Paper Roadmap}
\miniabstract{We outline the paper's contribution and workflow: a simulation-first replication of Avellaneda--Stoikov, followed by systematic stress-tests that vary only the mid-price model while keeping the execution and accounting engine fixed, and finally a covariance-aware multivariate extension. This isolates how price dynamics and dimensionality affect the naive-versus-inventory-based performance gap.}

Starting from the Avellaneda--Stoikov market making framework, this article uses Monte Carlo simulation to reproduce the canonical qualitative result: relative to a naive symmetric quoting rule, an inventory-based strategy improves risk-adjusted performance by actively controlling inventory exposure, typically reducing the dispersion of terminal wealth and increasing Sharpe-like metrics under comparable experimental conditions. The study is deliberately data-free: rather than calibrating to historical order book data, we adopt stylized parameter regimes and scenario grids, with the goal of isolating mechanisms and ensuring transparent reproducibility.

We begin with the baseline arithmetic Brownian motion (ABM) mid-price under a finite horizon, which provides the cleanest environment to interpret inventory risk and spread capture. We then replace ABM with geometric Brownian motion (GBM) to reflect non-negative prices and multiplicative fluctuations, and we discuss the practical consequences of price-level dependent risk when evaluating utility-based objectives and quoting policies. Finally, we rebuild the multivariate step as a covariance-aware inventory-control template: prices, inventories, fills, and quotes become vector-valued, and the scalar term $\sigma^2 q_t$ is replaced by the portfolio pressure $\Sigma q_t$.

\section*{Market making history}
This paper studies the interface between market making and market microstructure through the lens of two time scales. At the daily scale, an asset price can be idealized as a continuous-time process (e.g., a Brownian-motion-like mid-price), and the market maker's problem resembles dynamic risk management with the objective of harvesting the bid--ask spread while controlling inventory exposure. At the sub-second scale, the same environment becomes intrinsically discrete and rule-driven: prices live on a tick grid, orders arrive asynchronously, and the limit order book mechanics and trading frictions shape both the observed price formation and the feasibility of any quoting strategy. The core message is that market making cannot be fully understood at only one resolution: a smooth diffusion view is useful for tractability and control, while a microstructural view is essential for interpreting what the mid-price represents and for implementing strategies under real trading constraints.

Within this perspective, market making can be summarized as an algorithmic decision rule that continuously updates bid and ask quotes as a function of a state vector. Conceptually, the market maker forms a (possibly personal) estimate of value and quotes around it, posting a bid below and an ask above, then adapts these quotes as the state changes, most importantly as the mid-price moves and inventory accumulates. The economic role of this activity is liquidity provision: by standing ready to trade on both sides, the market maker supplies immediacy to other market participants who value execution without delay.

Historically, the intuition of a two-sided quote originated in floor-based settings with designated dealers (often specialists) who were central to trading in a given name. Modern electronic markets replace this institutional structure with competitive, automated quoting in data-center environments, where many agents simultaneously post and cancel orders at high frequency. This evolution motivates the academic progression from classical inventory-based dealer models, built around utility maximization and Hamilton--Jacobi--Bellman (HJB) equations, to electronic-market approximations that yield implementable quoting formulas and accommodate realistic constraints such as inventory bounds, risk limits, adverse selection, and multi-asset interactions.

Against this background, the Avellaneda--Stoikov paradigm provides a canonical benchmark: it formalizes the idea that the market maker effectively buys low and sells high, but on compressed horizons (seconds or minutes rather than days), and it makes explicit how optimal spreads and skew depend on inventory risk and order arrival dynamics. The goal of this paper is to apply this inventory-based logic across alternative stochastic specifications for the underlying state variables (different mid-price dynamics and a covariance-aware multivariate inventory state), while keeping the economic mechanism invariant: the market maker controls quotes to balance spread capture against inventory and execution risk under microstructural frictions.

\section*{Notation}
\begin{tabularx}{\textwidth}{@{}lX@{}}
\toprule
$S_t$ & Mid-price / reference price \\
$q_t$ & Inventory (number of shares/contracts) \\
$X_t$ & Cash account \\
$W_t$ & Total wealth / overall P\&L \\
$\delta_t^{bid}, \delta_t^{ask}$ & Distance from the mid-price for bid/ask quotes \\
$p_t^{bid}, p_t^{ask}$ & Quoted prices (bid/ask) \\
$\lambda^{bid}(\delta), \lambda^{ask}(\delta)$ & Arrival intensities of market orders hitting the bid/ask quotes \\
$N_t^{bid}, N_t^{ask}$ & Counting processes of fills at the bid/ask \\
$\sigma$ & Volatility / instantaneous scale \\
$\gamma$ & Risk aversion (CARA / exponential utility) \\
$T$ & Policy horizon \\
$\Delta t$ & Simulation time step \\
\bottomrule
\end{tabularx}

\chapter{Market Microstructure and High-Frequency Trading: LOB Foundations, Taxonomy, and Risks}
\label{ch:microstructure_hft}

\section{LOB Mechanics}
Modern electronic equity and crypto venues are typically organized as continuous double auctions implemented through a limit order book (LOB). At any time $t$, the book aggregates standing intentions to buy and sell. The best available buy price is the \emph{best bid} $p_t^{bid}$ and the best available sell price is the \emph{best ask} $p_t^{ask}$. The \emph{bid--ask spread} is
\[
\mathrm{spr}_t = p_t^{ask}-p_t^{bid},
\]
and a standard reference price is the \emph{mid-price}
\[
S_t=\frac{p_t^{ask}+p_t^{bid}}{2}.
\]
In practice, $S_t$ is an informational proxy rather than a traded price: executions occur at discrete prices on the book, while $S_t$ is a convenient state variable used to define quote distances $\delta_t^{bid},\delta_t^{ask}$ and to mark-to-market total wealth $W_t$ in reduced-form models.

A central microstructural friction is price discreteness. Exchanges impose a \emph{tick size} $\tau>0$, meaning that admissible prices lie on a grid (multiples of $\tau$). Tick size interacts with quoting and execution probability: for a given $S_t$, a one-tick change in $p_t^{bid}$ or $p_t^{ask}$ may cause a discontinuous jump in competitiveness and thus in expected fills.

The depth of the book describes how many shares/contracts are available at each price level. If $D_t^{bid}(p)$ denotes cumulative buy depth at prices $\ge p$ and $D_t^{ask}(p)$ cumulative sell depth at prices $\le p$, then the slope of these curves around the top of book summarizes local liquidity. Depth matters because it shapes the impact of market orders, the stability of the top quotes, and the typical waiting time for a new limit order to be executed.

Finally, most venues implement price-time priority. Orders at better prices execute first. Among orders at the same price, earlier arrivals have priority. This induces queueing: a market maker who posts at the best bid or ask typically joins a queue, and the probability of being filled depends not only on the level of $p_t^{bid}$ or $p_t^{ask}$ but also on queue position. These mechanics are key for realistic implementation, even when the analytical model abstracts them into a reduced-form \emph{arrival intensity} $\lambda^{bid}(\delta)$ and $\lambda^{ask}(\delta)$ that depends only on the distance $\delta$ from $S_t$. Standard treatments of these concepts can be found in classical market microstructure references such as \textcite{OHara1995}, \textcite{Harris2003}, and \textcite{Hasbrouck2007}.

\section{Orders and Order Flow}
LOB dynamics are driven by the interaction of different order types. A limit order specifies a side (buy/sell), a price, and a quantity, and provides liquidity by adding depth to the book. A market order consumes liquidity by executing immediately against the standing best quote(s), walking the book if its size exceeds the available top-of-book depth.

A key empirical feature is that books are highly dynamic because orders are frequently cancelled or modified. Cancellations are essential for market making: when $S_t$ moves or $q_t$ changes, a market maker typically cancels stale quotes and reposts new ones. From a modeling perspective, cancellations generate a high-frequency churn that affects queue position and realized execution rates; in reduced-form models, their net effect is often captured through effective \emph{fill processes} $N_t^{bid}$ and $N_t^{ask}$ with intensities $\lambda^{bid}(\cdot)$ and $\lambda^{ask}(\cdot)$.

Executions can also be partial. When a market order hits a limit order of size larger than the incoming quantity, only a portion is executed and the remaining quantity stays in the book, preserving price-time priority at that level. Conversely, if a market order size exceeds top depth, it consumes multiple levels and produces immediate changes in $p_t^{bid}$ and $p_t^{ask}$, which in turn shift $S_t$.

Latency is another structural element. There is always a non-zero delay between observing the state of the book and having a cancel/replace action acknowledged by the exchange. During this window, an HFT market maker is exposed to adverse price movements and unexpected executions at quotes that were intended to be updated. In practice, latency and message processing limits shape optimal quoting, but in baseline Avellaneda--Stoikov-type models they are typically abstracted away so that quotes can be treated as continuously adjustable. The resulting model should thus be interpreted as an idealized benchmark rather than a full description of the LOB.

These ingredients motivate the modeling choice adopted later: rather than simulating the entire queue and message flow, we represent executions by counting processes $N_t^{bid},N_t^{ask}$ whose conditional intensities $\lambda^{bid}(\delta)$ and $\lambda^{ask}(\delta)$ depend on how far the quotes are from the reference mid-price $S_t$. This captures the first-order trade-off faced by a liquidity provider: tighter $\delta_t^{bid},\delta_t^{ask}$ increase expected fills and spread capture, but also increase exposure to inventory risk through $q_t$ and to execution risk through random fill timing.

\section{Empirical Stylized Facts and Reduced-Form Book Models}
Empirical work on limit order books documents robust regularities across assets and venues. One such regularity is order book imbalance: when buy-side depth near the top of book exceeds sell-side depth (or vice versa), short-term price moves tend to be biased toward the thinner side. Imbalance provides a microstructural signal that can be used to tilt quoting or control inventory, although incorporating it explicitly requires either book data or a structural model for depth dynamics.

A second feature is resilience: after liquidity is consumed by aggressive orders, the book tends to replenish over time through new limit orders and cancellations. Resilience determines how quickly spreads revert and how persistent the impact of temporary shocks is. For market makers, low resilience increases the risk of being caught with inventory $q_t$ during sustained adverse moves in $S_t$, while high resilience improves the ability to mean-revert $q_t$ through subsequent fills.

A third feature is clustering in order flow. Trades and order submissions are not independent in time; instead, they arrive in bursts. This motivates models beyond homogeneous Poisson processes, such as self-exciting point processes, which can reproduce clustered arrivals and endogenous fluctuations in execution intensity. However, these richer specifications complicate both analysis and calibration.

Tick size discreteness creates additional effects. When $\tau$ is large relative to typical short-term volatility $\sigma_t$, the top of book may be relatively stable and market makers can earn the spread more reliably. When $\tau$ is small, competition tightens spreads and queueing becomes more relevant. In both cases, tick constraints imply that the mapping from a continuous reference $S_t$ to admissible quotes is quantized, which can be important when interpreting theoretical optimal $\delta_t^{bid},\delta_t^{ask}$ in implementation.

Comprehensive surveys of empirical LOB properties and modeling approaches are provided in \textcite{GouldEtAl2013}. In the present work, we intentionally adopt a reduced-form execution model for comparability across mid-price specifications: the execution mechanism (counting processes $N_t^{bid},N_t^{ask}$ with intensities $\lambda^{bid}(\cdot),\lambda^{ask}(\cdot)$) is held fixed while the mid-price dynamics are varied under a common finite-horizon design. This choice isolates the effect of the price model on inventory-based control, while acknowledging that richer microstructure features (imbalance signals, clustered flow, queue position, and latency) are natural extensions beyond the baseline benchmark.

\section{High-Frequency Trading: Operational Definitions and Key Metrics}
High-frequency trading (HFT) refers to algorithmic strategies that operate on short horizons, where the interaction between trading decisions and market microstructure becomes first-order. In this context, the \emph{policy horizon} $T$ is not only a modeling choice but an operational constraint: strategies are designed to make decisions and manage risk over seconds, milliseconds, or even microseconds, and their profitability depends on how quickly they can adapt to changing order flow and movements in the reference price $S_t$.

As seen in the previous section, a central metric is latency, broadly defined as the end-to-end delay between (i) observing a market state and (ii) having an action (e.g., a new quote, a cancel/replace, or a hedge) acknowledged and made effective by the venue. Latency can be decomposed into data latency, decision latency, and execution latency. In market making, latency directly impacts the ability to maintain intended quote distances $\delta_t^{bid},\delta_t^{ask}$ and inventory targets for $q_t$, and it creates exposure to fills at stale quotes, thereby affecting the realized dynamics of $N_t^{bid}$ and $N_t^{ask}$.

Turnover measures how frequently capital is recycled. At the strategy level, it is commonly proxied by traded volume relative to a notional base, and at the position level it is related to the rate at which $q_t$ mean-reverts. High turnover is not automatically desirable: it may increase gross spread capture, but it can also amplify transaction costs, queue losses, and exposure to toxic flow.

Because HFT operates at the micro scale, practitioners rely on micro-signals and liquidity indicators that summarize the state of the book and recent order flow. Typical examples include the spread $\mathrm{spr}_t$ and top-of-book depth, imbalance, recent trade signs and intensities, short-term realized volatility, and measures of flow persistence. These features are used either to tilt quoting (in market making), to time entry/exit (in arbitrage or stat-arb), or to control urgency and slippage (in execution). In this thesis, we keep the execution model reduced-form and let liquidity conditions primarily enter through the parameters governing $\lambda^{bid}(\delta),\lambda^{ask}(\delta)$ and through the volatility scale $\sigma$.

\section{Classes of Strategies}
HFT is best viewed as a family of strategies rather than a single paradigm. The common thread is fast reaction to market conditions, but objectives and risk exposures differ substantially across classes.

\paragraph{Market making (core of this work).}
A market maker supplies liquidity by posting bid and ask quotes and earns the spread conditional on being filled. The defining risk is inventory: if buys and sells do not balance, the market maker accumulates an inventory $q_t$ that is exposed to movements in $S_t$. In reduced-form models, the strategy controls the trade-off between execution probability and risk by adjusting quote distances $\delta_t^{bid},\delta_t^{ask}$ from the reference mid-price $S_t$. The present work treats market making as a controlled quoting problem, comparing a naive symmetric policy to an inventory-aware policy derived from the Avellaneda--Stoikov principle.

\paragraph{Statistical arbitrage and cross-venue trading (operational context).}
Statistical arbitrage encompasses strategies that exploit temporary mispricings across related instruments (pairs or baskets) or across venues. In a cross-venue setting, the strategy may buy on one venue and sell on another to capture price differentials, subject to latency, fees, and execution risk. In a related-instruments setting, the strategy may trade a spread or mean-reverting residual, requiring careful control of model risk and regime changes. These strategies share important infrastructure with market making (low-latency execution, queue management, and robust risk controls) but their alpha mechanism is typically driven by relative value rather than spread capture.

\paragraph{Execution and optimal execution (comparison baseline).}
Execution algorithms aim to trade a given quantity while minimizing cost and risk, rather than to provide liquidity continuously. The objective is typically a functional of expected slippage, volatility risk, and market impact, leading to policies that schedule participation rates or adaptively vary urgency over time. Classical frameworks formalize this problem using stochastic control and dynamic programming (see \textcite{CarteaJaimungalPenalva2015}). Execution provides a useful contrast to market making: in market making the trader chooses quotes to be hit, whereas in execution the trader chooses aggressiveness to achieve a target quantity, often paying the spread instead of earning it.

\section{Risks and Constraints}
Although HFT strategies can be framed as optimization problems, real-world performance is shaped by risks and constraints that may not be fully captured by baseline models. For market making, these risks motivate the inventory-aware structure studied later.

\paragraph{Inventory risk.}
Inventory risk is the exposure of the dealer's position to price movements. In the Avellaneda--Stoikov setting, inventory enters the objective through marked-to-market wealth $W_t$ and is controlled via quote skew: the market maker shifts $\delta_t^{bid}$ and $\delta_t^{ask}$ to encourage trades that reduce $q_t$. Inventory limits (hard bounds) and risk budgets (soft penalties) are standard constraints in practice.

\paragraph{Adverse selection and toxicity.}
Adverse selection occurs when incoming market orders are informational: the market maker sells just before $S_t$ rises or buys just before $S_t$ falls. This is often described as toxic flow, because it turns apparent spread revenue into losses. In reduced-form models, adverse selection can be approximated by state-dependent intensities $\lambda^{bid}(\cdot),\lambda^{ask}(\cdot)$ or by coupling executions to future changes in $S_t$, but fully modeling it typically requires richer microstructure inputs.

\paragraph{Queue risk.}
When quotes are placed at competitive prices, fills depend on queue position at $p_t^{bid}$ and $p_t^{ask}$. Frequent cancellations and reposting may improve alignment with $S_t$ but can worsen priority. Thus, there is a trade-off between staying ``fresh'' and maintaining queue position. Queue risk is particularly important in low-volatility regimes where many market makers crowd the top of book and marginal priority determines realized fill rates for $N_t^{bid}$ and $N_t^{ask}$.

\paragraph{Latency and implementation risk.}
Latency creates a gap between intended and realized control. During delays, the strategy may be executed at stale quotes, miss hedges, or experience rapid changes in $q_t$. Additional implementation constraints include exchange message limits, minimum resting times, order throttles, and fees/rebates structures that can materially alter optimal quoting behavior.

\paragraph{Regime shifts and model risk.}
HFT strategies are sensitive to regime changes in volatility, liquidity, and order flow. A regime shift can invalidate parameters that were stable in calm periods (e.g., the effective intensity scale in $\lambda^{bid}(\cdot),\lambda^{ask}(\cdot)$ or the volatility level $\sigma$), causing abrupt changes in fill rates and inventory dynamics.

\section{Bridge to Avellaneda--Stoikov}
The discussion above provides the minimum context needed to motivate the Avellaneda--Stoikov benchmark. The LOB supplies the two-sided quoting mechanism through $p_t^{bid},p_t^{ask}$ around a reference $S_t$, order flow and cancellations motivate a reduced-form execution model in terms of fill processes $N_t^{bid},N_t^{ask}$ and intensities $\lambda^{bid}(\cdot),\lambda^{ask}(\cdot)$ and HFT metrics and risks clarify why inventory-aware control of $q_t$ and robustness across regimes are central. With this foundation, we now introduce Avellaneda--Stoikov as a tractable stochastic control formulation of market making and use it as the baseline engine for the simulation-based comparisons developed in the remainder of the work.

\chapter{High-Frequency Trading in a Limit Order Book: Core Ideas Behind Avellaneda--Stoikov}
\label{sec:hft_lob_core_ideas}

\section{Problem statement: how does a market maker create a bid--ask?}
The central question is: \emph{how can a market maker (MM) generate bid and ask quotes in a limit order book?} The MM provides liquidity by continuously posting a two-sided market, aiming to profit from spread capture while controlling the risks induced by random price fluctuations and random execution timing.

Two main sources of uncertainty drive the problem. First, there is inventory risk, stemming from the stochastic nature of the reference (mid) price process: if the MM accumulates inventory $q_t$, the marked-to-market value of that inventory fluctuates randomly with $S_t$. Second, there is transaction risk, due to the randomness in both the number of market orders and their arrival times: executions occur at random instants, and the MM cannot perfectly control when (and how often) its quotes are hit. The goal is therefore to choose a quoting policy that maximizes the MM's objective, formulated as utility maximization.

\section{Two quoting paradigms: naive symmetric vs inventory-based}
A convenient way to organize the discussion is to contrast two approaches.

The first approach is a naive symmetric strategy. The MM posts a bid and an ask that are symmetric around the mid-price and do not depend on inventory. In other words, the distances from the mid-price are fixed (or exogenously time-dependent) and ignore the state variable $q_t$. This policy can capture spread revenue, but it does not actively control inventory risk: any inventory stabilization occurs only indirectly (e.g., through reduced fill rates when quoting wider).

The second approach is an inventory-based strategy: the MM calibrates its bid-ask quotes to the current inventory so as to push $q_t$ back toward $0$ by the end of the horizon. The key conceptual workflow is a two-step procedure. In Step~1, one computes an indifference price (or reservation price) that depends on inventory. In Step~2, one constructs bid and ask quotes around that inventory-dependent reference and calibrates their aggressiveness using the book/execution mechanism, typically summarized by arrival intensities.

\section{Minimal building blocks: price dynamics, order arrivals, and order sizes}
Even in a reduced-form model, the market-making problem requires three minimal blocks.

\paragraph{Mid-price dynamics.}
The reference price $S_t$ must be modeled, together with its volatility structure. Classical specifications use Brownian/It\^o diffusions, while richer specifications use Volterra-type dynamics or rough volatility drivers. Changing the dynamics of $S_t$ changes optimal quoting because it changes the distribution of future price moves and therefore the risk of holding inventory.

\paragraph{Order arrival mechanism.}
Executions are modeled via point processes (often Poisson in the baseline). A common reduced-form assumption is that the arrival intensity of market orders hitting the MM quotes depends on the quote distance from the mid-price. In notation consistent with the rest of the manuscript, one often writes effective intensities
\[
\lambda^{bid}(\delta), \qquad \lambda^{ask}(\delta),
\]
where $\delta$ represents the distance of the quote from $S_t$ (e.g., $\delta_t^{bid}$ or $\delta_t^{ask}$). The fill events are then summarized by counting processes $N_t^{bid}$ and $N_t^{ask}$.

\paragraph{Order size specification.}
One must specify the size of executions. In many baseline treatments, execution size is simplified to one unit per fill, which yields an inventory jump of $\pm 1$ at each execution. More general frameworks treat order sizes as random marks (marked point processes), which is more realistic but also more complex.

A practical remark is that under rough volatility or Volterra-type models, the volatility forecast over the relevant horizon $\sigma_t$ is no longer well-described by a constant (or even Markovian) but it is variable in the time. This changes how aggressively the MM must move quotes in response to inventory: when short-horizon volatility is clustered and persistent, inventory control often requires more state-dependent and regime-sensitive adjustments.

\section{A baseline mid-price model and CARA utility: why closed forms appear}
To illustrate how the indifference-price logic arises, consider a baseline specification where the mid-price follows an arithmetic Brownian motion with constant volatility. Specifically, assume that, for $u\ge t$,
\[
dS_u = \sigma\, dW_u,
\qquad S_t=s,
\]
where $W$ is a one-dimensional Brownian motion and $\sigma$ is constant. The modeling implication is that we impose no structure on drift or autocorrelation beyond this diffusion assumption. In particular, there is no attempt to encode predictability in $S_t$.

The objective is to maximize the exponential utility of terminal P\&L at a finite horizon $T$, i.e., a CARA (constant absolute risk aversion) preference with risk-aversion parameter $\gamma>0$. Exponential utility also connects naturally to convex risk measures and provides a clean framework for defining indifference prices.

As a first step, consider a frozen-inventory thought experiment in which the inventory is held fixed at level $q$ between $t$ and $T$. The associated value is
\[
v(x,s,q,t)
=
\mathbb{E}_t\!\left[-\exp\!\left(-\gamma\bigl(x+qS_T\bigr)\right)\right],
\]
where $x$ denotes initial wealth at time $t$, $s$ is the current mid-price $S_t$, and $q$ is the number of shares held. The parameter $\gamma$ quantifies risk aversion: increasing $\gamma$ makes the agent more risk averse (larger penalty for variance), while decreasing $\gamma$ moves the agent closer to risk neutrality.

Under arithmetic Brownian motion,
\[
S_T = s + \sigma\bigl(W_T-W_t\bigr),
\]
and conditioning on $\mathcal{F}_t$ yields
\[
\mathbb{E}_t[S_T]=s,
\qquad
\mathrm{Var}_t(S_T)=\sigma^2(T-t).
\]
Define terminal wealth
\[
Y:=x+qS_T.
\]
Then
\[
\mathbb{E}_t[Y]=x+qs,
\qquad
\mathrm{Var}_t(Y)=q^2\sigma^2(T-t).
\]
Using the standard Gaussian moment generating identity, if $Y\sim N(m,v)$ then 
\begin{align*}
    \mathbb{E}(e^{aY})=\exp(am+\tfrac12a^2v)
\end{align*}
with $a=-\gamma$, one obtains
\[
\mathbb{E}_t\!\left[-e^{-\gamma Y}\right]
=
-\exp\!\left(
-\gamma\,\mathbb{E}_t[Y]
+\frac12\gamma^2\,\mathrm{Var}_t(Y)
\right)
=
-\exp\!\left(
-\gamma(x+qs)
+\frac12\gamma^2 q^2\sigma^2(T-t)
\right).
\]
Equivalently, by splitting exponentials,
\[
v(x,s,q,t)
=
-\exp(-\gamma x)\,
\exp(-\gamma q s)\,
\exp\!\left(\frac{\gamma^2 q^2 \sigma^2 (T-t)}{2}\right).
\]
This factorization is particularly clean because the terminal wealth is Gaussian under ABM and the exponential utility interacts well with Gaussian mgfs. If instead $S_t$ followed GBM, such a neat factorization generally does not appear.

This ``closed-form friendliness'' is a direct consequence of the CARA/exponential specification $U(w)=-e^{-\gamma w}$: under normality (or more generally whenever the relevant moment generating function exists in a tractable form) computations simplify substantially. In Black--Scholes, for instance, $S_t$ is lognormal and expectations of the form $\mathbb{E}\bigl[-e^{-\gamma q S_T}\bigr]$ can be interpreted via Laplace-transform-type objects. For rough or otherwise more complex models, closed forms are typically lost and one must either (i) formulate and solve the associated Hamilton--Jacobi--Bellman (HJB) equation when the control is continuous, (ii) rely on numerical methods/simulation to approximate the value function, or (iii) modify the objective (e.g., mean--variance criteria or utilities on transformed prices such as log-prices).

Importantly, in more complex, non-Gaussian diffusion models, closed forms are typically lost but the qualitative dependence between optimal bid--ask placement and inventory remains. What changes is the computation: in Gaussian settings one obtains closed forms and ``clean'' separability, while in more realistic settings the inventory correction typically becomes a function of the full state (e.g., $s,q,t,\text{volatility},\text{other states}$) because both risk and the future distribution of $S_t$ depend on those factors. Operationally, the optimality conditions may be expressed in terms of partial derivatives of the value function, which are difficult to obtain in high-dimensional HJB/PDE problems and therefore motivate numerical schemes.

\section{Reservation prices via indifference: finite horizon}
The key quantity in inventory-based quoting is the \emph{reservation price} (indifference price): the price that makes the agent indifferent between keeping the current portfolio and switching to a portfolio with one more (or one fewer) share.

Formally, define the reservation bid price $r^b(s,q,t)$ and reservation ask price $r^a(s,q,t)$ as the solutions to
\[
v\bigl(x-r^b(s,q,t),\,s,\,q+1,\,t\bigr)=v(x,s,q,t),
\qquad
v\bigl(x-r^a(s,q,t),\,s,\,q-1,\,t\bigr)=v(x,s,q,t).
\]
Under the ABM--CARA frozen-inventory expression above, these conditions yield
\[
r^b(s,q,t)
=
s+(-1-2q)\,\frac{\gamma\sigma^2(T-t)}{2},
\qquad
r^a(s,q,t)
=
s+(1-2q)\,\frac{\gamma\sigma^2(T-t)}{2}.
\]
A convenient reference is the ``center'' reservation level (the midpoint between reservation bid and ask),
\[
r(s,q,t)
=
\frac12\bigl(r^a(s,q,t)+r^b(s,q,t)\bigr)
=
s-\frac12\,\gamma q\sigma^2(T-t).
\]
This expression highlights the inventory effect: if the agent is long stock ($q>0$), then $r(s,q,t)$ lies below the mid-price, meaning the agent is willing to sell at a lower price to reduce inventory. If the agent is short ($q<0$), then $r(s,q,t)$ lies above the mid-price, meaning the agent is willing to buy back at a higher price to close the short position.

In this paper, the reservation-price construction is kept in its finite-horizon form. This is sufficient for the ABM and GBM simulation benchmarks and keeps the terminal-wealth criterion consistent across the single-asset models.

The reservation-price step above provides the first half of the Avellaneda--Stoikov logic. The second half connects reservation prices to actual bid-ask quotes by incorporating execution probabilities through $\lambda^{bid}(\cdot)$ and $\lambda^{ask}(\cdot)$, i.e., by choosing quote distances that trade off fill rates against inventory and price risk. This connection is developed in the subsequent Avellaneda--Stoikov section.

\section{Market Making as a Liquidity Provision Problem}
A market maker provides immediacy by standing ready to buy and sell. If the dealer posts a bid below the current reference price and an ask above it, profits arise when both sides are hit over time and the spread is captured. However, the execution of the two sides is random. If the dealer buys repeatedly without selling, inventory builds up; if the price then falls, marked-to-market wealth deteriorates. Inventory control is therefore central to market making, especially on short horizons where positions can grow quickly relative to the dealer's risk budget.

Historically, two-sided quoting was associated with designated specialists or dealers on centralized exchanges. In modern electronic markets, the same economic function is carried out by many competing agents that post and cancel orders at high frequency. This institutional evolution is why reduced-form models of quoting, execution intensities, and inventory risk remain valuable: they transform a complex market mechanism into a control problem with clear state variables and interpretable parameters.

\chapter{The Avellaneda--Stoikov Framework}
\miniabstract{This chapter introduces the state variables, the reduced-form order-arrival mechanism, and the utility-based objective that together define the baseline market making problem.}

\section{State Dynamics}
The state of the MM is fully characterized by the triad $(S_t, q_t, X_t)$, representing respectively the reference mid-price, the current inventory (number of shares/contracts), and the cash account. The dealer provides liquidity by posting quotes at distances $\delta_t^{bid}$ and $\delta_t^{ask}$ from the mid-price
\[
p_t^{bid}=S_t-\delta_t^{bid}, \qquad p_t^{ask}=S_t+\delta_t^{ask},
\]
with $\delta_t^{bid},\delta_t^{ask}\ge 0$.

It is crucial to distinguish the continuous nature of the price from the discrete nature of the dealer's accounting. Assuming unit trade sizes, the executions are modeled via the counting point processes $N_t^{bid}$ and $N_t^{ask}$. Consequently, inventory and cash are jump processes that evolve only when an order is filled
\begin{align*}
    &dq_t=dN_t^{bid}-dN_t^{ask},\\
    &dX_t=(S_t+\delta_t^{ask})\,dN_t^{ask}-(S_t-\delta_t^{bid})\,dN_t^{bid}.
\end{align*}
Economically, when a market seller hits the bid ($dN_t^{bid}=1$), the dealer buys one share (increasing $q_t$) and pays the bid price (decreasing $X_t$). Conversely, a fill at the ask ($dN_t^{ask}=1$) reduces inventory and increases cash.

The continuous performance of the strategy is measured by the total wealth at any time $t$
\[
W_t=X_t+q_tS_t.
\]
This equation encapsulates the core of the market making problem. While the accumulated cash $X_t$ is insulated from price movements, the portfolio value $q_tS_t$ is continuously exposed to the stochastic fluctuations of $S_t$. The terminal objective evaluated at the policy horizon $T$ is based on 
\[
W_T=X_T+q_TS_T,
\]
driving the need for an optimal control policy that balances the spread capture in $X_t$ against the inventory risk inherent in $q_tS_t$.

\section{Arrival Intensities}
To close the specification of the state dynamics, we must define the behavior of the counting processes $N_t^{bid}$ and $N_t^{ask}$. The dealer does not have direct control over when an execution occurs, but they control the probability of an execution by choosing the quotation distances $\delta_t^{bid}$ and $\delta_t^{bid}$.

Let $\lambda(\delta)$ denote the arrival rate (intensity) of filled limit orders at a distance $\delta$ from the mid-price. Following the specific framework proposed by Avellaneda and Stoikov, we assume that the arrival of executions depends on three main market components:
\begin{enumerate}
    \item The overall constant frequency of market orders (liquidity-taking orders), denoted by $\Lambda$;
    \item The distribution of the volume size $Q$ of these market orders, which is empirically shown to follow a power law density $f^Q(x) \propto x^{-1-\alpha}$;
    \item The temporary market impact $\Delta p$ caused by a large market order.
\end{enumerate}
If a market order creates a temporary price impact $\Delta p$, the dealer's limit order posted at a distance $\delta$ will be executed only if the impact is large enough to reach the quote, i.e., if $\Delta p>\delta$. Therefore, the execution intensity is proportional to the probability of this event.

Following the finding that the price impact of an order of size $Q$ is best described by a logarithmic function 
\[
\Delta p \propto \ln{(Q)} 
\]
we can derive the exact functional form for the arrival intensity as follows
\begin{align*}
    \lambda(\delta)&=\Lambda \mathbb{P}(\Delta p>\delta)=\\
    &=\Lambda \mathbb{P}(\ln{(Q)}>K\delta)=\\
    &=\Lambda \mathbb{P}(Q>e^{K\delta})=\\
    &=\Lambda \int_{e^{K\delta}}^{+\infty}x^{-1-\alpha}dx= Ae^{-k\delta},
\end{align*}
where the constants are redefined as $A = \Lambda /\alpha$ and $k=\alpha K$.

This result provides us with the exponential intensities for the bid and ask sides
\[
\lambda^{ask}(\delta)=Ae^{-k\delta}, \qquad \lambda^{bid}(\delta)=Ae^{-k\delta}.
\]
In this reduced-form representation, the parameter $A$ retains the interpretation of the base market activity (the arrival rate at the mid-price), while $k$ controls execution sensitivity to quote aggressiveness. A larger $k$ implies a thicker book where the probability of execution decays rapidly as the dealer widens the spread. This dynamic establishes the core trade-off of the market making problem: quoting further away secures a larger margin per trade but exponentially dampens the execution probability.

\section{Objective Function}
Having established the dynamics of the state variables and the microstructural mechanics of order executions, we now formalize the dealer's objective. The market maker seeks to maximize the expected utility of their terminal wealth at a fixed horizon $T$ (e.g., the end of the daily trading session). Let 
\[
W_T = X_T + q_TS_T
\]
denote the marked-to-market wealth at the terminal time $T$. To rigorously capture the dealer's aversion to inventory risk, the optimization is cast within an expected utility framework using a Constant Absolute Risk Aversion (CARA) utility function
\[
u(W_T) = -e^{-\gamma W_T} 
\]
where $\gamma>0$ is the absolute risk aversion coefficient. The choice of the exponential utility is not merely mathematically convenient, it economically reflects the dealer's penalization of wealth variance. Since the cash account is riskless, the variance of the terminal wealth is entirely driven by the inventory position $q_T$ exposed to the price volatility $\sigma$.

The market maker's problem consists of finding the optimal predictable quoting policies (the controls), $\delta_t^{bid}$ and $\delta_t^{ask}$, that maximize the expected terminal utility conditional on the state at time $t$. We define the value function $u(t,S_t ,q_t, X_t)$ as
\begin{align*}
    u(t,s,q,x)&=\max_{\{\delta_t^{bid},\delta_t^{ask}\}_{t\in[0,T]}} \mathbb{E}\left[-\exp\!\left(-\gamma W_T\right)|S_t=s,q_t=q,X_t=x\right]=\\
    &=\max_{\{\delta_t^{bid},\delta_t^{ask}\}_{t\in[0,T]}} \mathbb{E}_t\left[-\exp\!\left(-\gamma (X_T+q_TS_T)\right)\right].
\end{align*}
This formulation represents a classic finite-horizon stochastic optimal control problem. Exponential utility is convenient because it interacts well with Gaussian benchmark models and yields interpretable reservation-price formulas \cite{AvellanedaStoikov2008,GueantLehalleFernandezTapia2013}.

Having formally defined the stochastic control problem and demonstrated the wealth-independence of the quoting strategy, the next step is to derive the optimal controls $\delta_t^{bid}$ and $\delta_t^{ask}$. This requires solving the dynamic programming problem through its associated Hamilton-Jacobi-Bellman (HJB) equation, which we establish in Chapter 4.

\chapter{Solving the Optimal Control Problem}
\miniabstract{The experimental comparison in every scenario is between a naive symmetric quoting rule and an inventory-based quoting rule derived from the Avellaneda--Stoikov control principle.}

\section{The Hamilton-Jacobi-Bellman Equation and Optimal Spreads}
As established in the previous chapter, the market maker's goal is to find the optimal quoting policies $\delta^{bid}$ and $\delta^{bid}$ that maximize the expected CARA utility of terminal wealth. According to the dynamic programming principle, this stochastic optimal control problem can be characterized by its infinitesimal generator. By applying Itô's Lemma for jump-diffusion processes, the problem translates into the Hamilton-Jacobi-Bellman (HJB) equation.

Assuming for now a generic Markovian dynamics for the mid-price $S_t$ characterized by an instantaneous variance $\sigma^2(t,S_t)=\sigma^2$, the HJB equation for the dealer's value function $u(t,s,q,x)$ takes the form

\begin{align*}
    \frac{\partial u}{\partial t} + \frac{1}{2}\sigma^2 \frac{\partial^2 u}{\partial s^2} &+ \max_{\delta^{bid}} \left\{ \lambda(\delta^{bid}) [u(t, s, q+1, x - s + \delta^{bid}) - u(t,s,q,x)] \right\} \\
    &+ \max_{\delta^{ask}} \left\{ \lambda(\delta^{ask}) [u(t, s, q-1, x + s + \delta^{ask}) - u(t,s,q,x)] \right\} = 0.
\end{align*}

This non-linear partial differential equation expresses a continuous equilibrium condition: the expected loss of utility due to time decay and price risk (captured by the continuous derivatives) must be exactly compensated by the expected gain from capturing the spread when limit orders are executed (the jump terms driven by the $\max$ operators). 

A critical feature of the exponential utility function $U(W) = -e^{-\gamma W}$ is that it allows the value function to be wealth-independent. We can factor out the cash component $x$ by applying
\begin{align*}
    u(t,s,q,x) = -e^{-\gamma x}e^{-\gamma \theta(s,q,t)},
\end{align*}
where $\theta(s,q,t)$ represents the certainty equivalent of the dealer's inventory position. Substituting this into the HJB equation and dividing by the non-zero common factor $u$, the cash variable $x$ is eliminated, and the problem reduces to an equation exclusively in terms of $\theta$
\begin{align*}
    \frac{\partial \theta}{\partial t} + \frac{1}{2}\sigma^2 \frac{\partial^2 \theta}{\partial s^2} - \frac{1}{2}\gamma \sigma^2 \left(\frac{\partial \theta}{\partial s}\right)^2 &+ \max_{\delta^{bid}} \frac{\lambda(\delta^{bid})}{\gamma} \left[ 1 - e^{-\gamma (s - \delta^{bid} - \theta(q) + \theta(q+1))} \right] \\
    &+ \max_{\delta^{ask}} \frac{\lambda(\delta^{ask})}{\gamma} \left[ 1 - e^{-\gamma (\theta(q-1) - \theta(q) - s - \delta^{ask})} \right] = 0.
\end{align*}
A powerful structural property of this formulation is that the optimization of the quoting distances can be decoupled from the explicit solution of the certainty equivalent $\theta$. To see this, we define the dealer's reservation prices for buying and selling an additional unit of inventory as the marginal change in the certainty equivalent
\begin{align*}
    &r^{bid}(s,q,t) = \theta(s,q+1,t) - \theta(s,q,t),\\
    &r^{ask}(s,q,t) = \theta(s,q,t) - \theta(s,q-1,t).
\end{align*}
Focusing on the bid side, the embedded optimization problem becomes 
\begin{align*}
    \max_{\delta^{bid}} \lambda(\delta^{bid}) [1 - e^{-\gamma(s - \delta^{bid} - r^{bid})}].
\end{align*}
The first-order condition with respect to $\delta^{bid}$ yields an exact, model-free relationship between the optimal spread and the elasticity of the limit order book
\begin{align*}
    s - r^{bid} =\delta^{bid} - \frac{1}{\gamma} \ln \left( 1 - \gamma \frac{\lambda(\delta^{bid})}{\lambda'(\delta^{bid})} \right).
\end{align*}
Under our assumption of exponentially decaying arrival intensities, $\lambda(\delta) = A e^{-k\delta}$, the ratio $\lambda / \lambda'$ simplifies strictly to $-1/k$. Therefore, the optimal quoting distances relative to the indifference prices are derived as
\begin{align*}
    &\delta^{bid} = (s - r^{bid}) + \frac{1}{\gamma} \ln \left( 1 + \frac{\gamma}{k} \right),\\
    &\delta^{ask} = (r^{ask} - s) + \frac{1}{\gamma} \ln \left( 1 + \frac{\gamma}{k} \right).
\end{align*}
This result is profoundly elegant: it demonstrates that the optimal width of the spread is determined primarily by the liquidity density of the order book $k$, independent of the underlying price dynamics. The compensation for the dealer's risk aversion and inventory position is, instead, entirely absorbed by the reservation prices $r^{bid}$ and $r^{ask}$, which act as the dynamic anchor for the quotes.

\section{The Certainty Equivalent and Reservation Price}
To make the optimal quoting distances derived in the previous section operational, we must explicitly determine the certainty equivalent $\theta(s,q,t)$ and, consequently, the reservation prices $r^{bid}$ and $r^{ask}$. 

As shown in the reduced HJB equation, the dynamics of $\theta$ depend heavily on the continuous diffusive term, which is driven by the underlying mid-price process. Following the Avellaneda-Stoikov baseline benchmark, we assume the mid-price follows an Arithmetic Brownian Motion (ABM) with constant volatility $\sigma$
\begin{align*}
    dS_t = \sigma dW_t.
\end{align*}
Under this specification, the pure inventory risk penalty can be isolated by solving the HJB equation for a "frozen inventory" portfolio (i.e., ignoring the jump terms associated with spread capture). In this case, the certainty equivalent admits a closed-form solution. The subjective value to the dealer of holding $q$ shares is simply the marked-to-market value of the portfolio minus a variance penalty
\begin{align*}
    \theta(s,q,t) = s \cdot q - \frac{1}{2}\gamma\sigma^2 q^2 (T-t).
\end{align*}
Applying the marginal difference definitions established in Section 4.1, we can compute the exact reservation bid and ask prices. The maximum price the dealer is willing to pay to acquire an additional share is
\begin{align*}
    r^{bid}(s,q,t) = \theta(s,q+1,t) - \theta(s,q,t) = s - q\gamma\sigma^2(T-t) - \frac{1}{2}\gamma\sigma^2(T-t).
\end{align*}
Conversely, the minimum price the dealer is willing to accept to sell a share is
\begin{align*}
    r^{ask}(s,q,t) = \theta(s,q,t) - \theta(s,q-1,t) = s - q\gamma\sigma^2(T-t) + \frac{1}{2}\gamma\sigma^2(T-t).
\end{align*}
For practical implementation, it is convenient to define a single, central reservation price $r(s,q,t)$ as the midpoint between the bid and ask reservation levels:
\begin{align*}
    r(s,q,t) = \frac{r^{bid}(s,q,t) + r^{ask}(s,q,t)}{2} = s - q\gamma\sigma^2(T-t).
\end{align*}

This formula represents the dealer's personalized ``fair value'' for the asset and constitutes the core of the inventory control mechanism. It clearly highlights the linear inventory effect under the ABM assumption: if the dealer is long ($q > 0$), the reservation price is shifted below the market mid-price, reflecting an urgency to sell. If short ($q < 0$), the reference price is shifted upwards. 

Crucially, the magnitude of this shift (and thus the strength of the inventory control) is directly proportional to the accumulated inventory $q$, the risk aversion parameter $\gamma$, the constant variance $\sigma^2$, and the time remaining to the end of the trading horizon $(T-t)$.

\section{The Inventory-Based Strategy (Avellaneda-Stoikov)}
By combining the results derived in the previous sections, we can finally define the operational market making strategy proposed by Avellaneda and Stoikov. The strategy is characterized by two distinct components: a microstructural component $\delta$, which dictates the width of the spread, and a risk-management component $r$, which dictates its position.

The optimal bid and ask quotes, $p^{bid}$ and $p^{ask}$, are obtained by centering the optimal distance $\delta$ around the central reservation price $r(s,q,t)$
\begin{align*}
    p^{bid}(s,q,t) = r(s,q,t) - \frac{1}{\gamma} \ln \left( 1 + \frac{\gamma}{k} \right),\\
    p^{ask}(s,q,t) = r(s,q,t) + \frac{1}{\gamma} \ln \left( 1 + \frac{\gamma}{k} \right).
\end{align*}
Substituting the closed-form expression for $r$ under the ABM assumption, the strategy can be explicitly written as
\begin{align*}
    &p^{bid} = \underbrace{s - q\gamma\sigma^2(T-t)}_{\text{Reservation Price}} - \underbrace{\frac{1}{\gamma} \ln \left( 1 + \frac{\gamma}{k} \right)}_{\text{Spread Distance } \delta},\\
    &p^{ask} = \underbrace{s - q\gamma\sigma^2(T-t)}_{\text{Reservation Price}} + \underbrace{\frac{1}{\gamma} \ln \left( 1 + \frac{\gamma}{k} \right)}_{\text{Spread Distance } \delta}.
\end{align*}

This formulation reveals the core logic of inventory-based market making: the dealer maintains a constant total spread, but its center is "skewed" relative to the market mid-price $s$ based on the current inventory $q$. 

Specifically, we can identify different behaviors based on the $q$ value. If $q=0$ (neutral inventory), the reservation price equals the mid-price ($r=s$). The dealer quotes symmetrically around the market price, acting as a pure liquidity provider. If $q>0$ (long position), the reservation price shifts downward ($r < s$). Both the bid and ask quotes are lowered and this makes the dealer's ask more competitive (increasing the probability of selling) and the bid less attractive (decreasing the probability of buying more shares), effectively inducing a mean-reverting behavior on the inventory. If $q<0$ (short position), the reservation price shifts upward ($r > s$). The quotes are raised, making the bid more competitive to attract sellers and the ask less attractive to avoid further shorting.

The intensity of this skewing effect is dynamically adjusted according to the remaining time horizon $(T-t)$. As the trading session approaches its end, the penalty for holding a non-zero inventory increases, forcing the dealer to be more aggressive in their rebalancing efforts. This sophisticated balance between market-neutral profit harvesting and active risk mitigation constitutes the primary advantage of the Avellaneda-Stoikov framework over naive quoting strategies.

\section{Benchmark Strategies for Simulation}

To systematically evaluate the performance and the risk-mitigation capabilities of the inventory-based policy derived in Section 4.3, we must establish a clear baseline for comparison. In the simulation experiments conducted in the subsequent chapters, the Avellaneda-Stoikov model will be tested against a primary benchmark: the Naive Symmetric Strategy.

The naive symmetric strategy represents the behavior of a traditional liquidity provider who optimizes for spread capture based on market microstructure, but completely ignores inventory risk. Mathematically, this behavior corresponds to a dealer who calculates the optimal quoting distance $\delta$ using the limit order book elasticity $k$, but blindly anchors these quotes to the current market mid-price $s$ rather than their personal reservation price $r$.

Using the same optimal distance $\delta$ defined previously, the quotes for the symmetric dealer are simply
\begin{align*}
    &p^{bid}_{sym}(s,t) = s - \delta = s - \frac{1}{\gamma} \ln \left( 1 + \frac{\gamma}{k} \right),\\
    &p^{ask}_{sym}(s,t) = s + \delta = s + \frac{1}{\gamma} \ln \left( 1 + \frac{\gamma}{k} \right).
\end{align*}
By keeping the quotes perfectly centered around the mid-price ($p^{bid}_{sym}$ and $p^{ask}_{sym}$ are equidistant from $s$), this strategy maximizes the probability of symmetric executions in a static market. However, because the reference price does not adjust to the accumulated inventory $q$, the symmetric dealer lacks any endogenous mechanism to mean-revert their position. 

If the market exhibits a strong directional trend, the symmetric dealer will continuously execute on one side of the book, accumulating a massive directional inventory (e.g., buying in a falling market). Without the "skewing" mechanism to protect the portfolio, the marked-to-market wealth $W_t$ of the symmetric strategy is exposed to severe drawdowns. 

Comparing the inventory-based strategy (which actively shifts $r$) against this symmetric benchmark (which keeps $r=s$) isolates the exact value of the inventory control mechanism. This dynamic will be the central focus of the event-driven simulations developed in the following chapters, where both strategies will be subjected to identical order flows and price paths across different diffusion assumptions.

\chapter{Classical Diffusions under a Finite Horizon}
\miniabstract{This chapter contains the core single-asset contribution of the paper. The same execution and inventory engine is tested under two finite-horizon benchmark cases to understand the role of the diffusion model.}

While Chapter 4 establishes the exact derivation for the baseline case, the subsequent scenarios introduce analytical complexities, such as the lack of a closed-form solution for the Geometric Brownian Motion (GBM), that in practice are addressed through heuristic approximations. Therefore, from this point forward, we will focus on the economic intuition of the reservation price and its numerical implementation within the simulation engine.

\section{Case 1: Arithmetic Brownian Motion with Finite Horizon}

In this initial benchmark scenario, we model the dynamics of the asset's mid-price through an Arithmetic Brownian Motion (ABM). Assuming a constant drift $\mu$ and a constant volatility parameter $\sigma$, the stochastic differential equation (SDE) governing the price process is defined as
\begin{align*}
    dS_t = \mu dt + \sigma dW_t.
\end{align*}
The decision to employ an ABM, deviating from the Geometric Brownian Motion (GBM) typically used in derivative pricing, is motivated by several structural and economic considerations. Primarily, the ABM provides the most transparent setting to isolate and interpret inventory risk, since price increments are additive, the marked-to-market effect of holding a position is strictly linear with respect to the inventory size. Furthermore, this price dynamics offers superior analytical tractability when paired with the Constant Absolute Risk Aversion (CARA) utility function. Under the ABM assumption, terminal wealth is normally distributed, allowing the moment-generating function to factorize and yielding elegant closed-form solutions for the dealer's reservation price. This choice also ensures that the utility functionals remain bounded throughout the optimization process.

The optimization in this case is conducted over a fixed time horizon $[0, T]$, which represents a standard trading session. Accounting for the drift term, the dealer's central reservation price takes the explicit form 
\begin{align*}
    r(s, q, t) = s + \mu(T-t) - q\gamma\sigma^2(T-t).
\end{align*}
This formulation highlights two fundamental temporal dynamics: the directional anticipation and the horizon effect. While the $\mu(T-t)$ term adjusts the valuation based on the expected short-term trend, the risk premium associated with the inventory ($q\gamma\sigma^2(T-t)$) decays linearly as the terminal time $T$ approaches. Intuitively, as the time-to-horizon $(T-t)$ diminishes, there is a lower probability that the mid-price will undergo significant adverse fluctuations before the final liquidation. Consequently, as $t$ approaches $T$, the indifference price $r(s, q, t)$ converges toward the mid-price $s$, and the dealer's quotes naturally realign with those of a symmetric strategy.

The primary objective of this experiment is to empirically measure the extent to which the inventory-aware strategy succeeds in reducing the dispersion of terminal wealth compared to a naive symmetric benchmark. By maintaining a constant order flow and fixed accounting metrics, the simulation aims to demonstrate that active quote skewing stabilizes profits by drastically reducing P\&L variance. Furthermore, we expect to observe an improvement in risk-adjusted returns, as evidenced by higher Sharpe ratios, and a more disciplined control over inventory excursions. The strategy should consistently drive the position $q_t$ toward neutrality before the end of the trading horizon, effectively mitigating the dealer's exposure to directional price risk.

\section{Case 2: Geometric Brownian Motion with Finite Horizon}

The second benchmark replaces the Arithmetic Brownian Motion with a Geometric Brownian Motion (GBM). In this setting, the dynamics of the mid-price are governed by the following stochastic differential equation
\begin{align*}
    dS_t = \mu S_t dt + \sigma S_t dW_t.
\end{align*}
This specification addresses a primary limitation of the ABM by preserving the strict positivity of asset prices. More importantly, the transition to GBM introduces a multiplicative price risk structure where the instantaneous volatility is proportional to the current price level $S_t$. Consequently, the inventory risk becomes price-level dependent: the absolute variance of a position scales quadratically with the asset price, meaning that holding a fixed number of shares $q$ at a higher price exposes the dealer to significantly greater marked-to-market risk.

Due to this non-linear structure, the dealer's optimization problem under CARA utility does not yield a simple, exact closed-form solution as in the ABM case. To overcome this, we implement a local variance approximation. By treating the price dynamics as locally arithmetic at each time step $t$, we define an effective volatility $\sigma_{eff} = \sigma S_t$, which represents the instantaneous monetary fluctuations of the asset. By substituting this state-dependent variance into the Avellaneda-Stoikov framework, we derive the following heuristic for the reservation price
\begin{align*}
    r(s, q, t) = s + \mu s(T-t) - q\gamma(\sigma s)^2(T-t).
\end{align*}
This formula, utilized within our simulation engine, allows the agent to dynamically scale its risk premium. As the price $s$ increases, the term $q\gamma(\sigma s)^2(T-t)$ grows quadratically, forcing the market maker to skew quotes more aggressively to induce inventory mean-reversion. 

The primary objective of this simulation is to observe how this adaptive quoting behavior performs compared to a naive strategy that fails to account for the multiplicative nature of the risk. We expect the inventory-aware agent to maintain tighter control over its position limits during price rallies, effectively mitigating the "exploding" variance inherent in the GBM process. This case serves as a crucial bridge between the idealized linear world of the ABM and the more realistic, heteroskedastic dynamics of actual financial markets.

\section{Synthesis: ABM vs GBM}

Having defined both the Arithmetic (ABM) and Geometric (GBM) price dynamics within a finite-horizon framework, it is crucial to synthesize the theoretical divergence between the two models before examining the empirical results of our event-driven simulation. Although both models share the underlying mechanics of order arrival and dealer accounting, the nature of the stochastic price path fundamentally alters the risk landscape for the market maker.

The divergence is rooted in the structure of the variance. In the ABM framework (Section 5.1), the inventory risk premium is strictly linear and price-independent: holding $q$ shares carries the same marked-to-market risk regardless of whether the asset trades at $\$10$ or $\$100$. This idealized, homoskedastic environment provides a clean baseline to observe the pure mechanics of the Avellaneda-Stoikov inventory control and the linear decay of the horizon effect as time $t$ approaches terminal time $T$.

Conversely, the transition to the GBM framework (Section 5.2) introduces heteroskedasticity into the price process. The instantaneous variance of the position scales quadratically with the asset's price level ($S_t^2$). Consequently, a dealer operating under GBM cannot rely on a static perception of risk. As the asset appreciates in value, the monetary impact of price fluctuations magnifies, demanding a highly adaptive quoting response. To mitigate this exploding variance, the inventory-aware agent utilizes a local variance approximation, substituting a constant $\sigma^2$ with an effective, state-dependent variance $(\sigma S_t)^2$.

However, despite this profound divergence in how inventory risk is priced into the reservation price, a crucial structural similitude persists between the two frameworks: the formulation of the optimal spread. The Avellaneda-Stoikov architecture gracefully separates the management of inventory risk from the optimization of execution probability. Consequently, the optimal distance between the dealer's quotes and the reservation price (the spread offset) remains mathematically invariant across both dynamics. This offset is strictly a function of the market's liquidity profile ($\kappa$) and the dealer's risk aversion parameter ($\gamma$), representing a pure liquidity premium. Therefore, while the GBM market maker aggressively shifts their entire quoting window to defend against price-scaled variance, the intrinsic width of that window remains identical to that of the ABM dealer.

\chapter{Simulation and Experimental Design}
\miniabstract{All model cases are studied through the same simulation engine. This isolates the effect of the price process from the execution mechanism itself.}

\section{The Time-Discretized Simulation Engine}
To numerically evaluate the strategies' performance, a custom simulation engine was implemented in Python. Contrary to a purely event-driven approach, which would be analytically complex when coupling continuous diffusion processes (mid-price) with discrete jump processes (market orders), the simulator relies on an Euler-Maruyama discretization scheme.

The time horizon $T$ is partitioned into $N$ equidistant intervals of length $\Delta t$. At each discrete time step $t_i$, given the state vector comprising the mid-price $S_{t_i}$, the accumulated inventory $q_{t_i}$, and the cash account $X_{t_i}$, the simulator executes the following logical cycle.

\begin{enumerate}
    \item \textbf{Quote Generation:} The active policy computes the optimal quote prices $p^{bid}$ and $p^{ask}$ based on the current state, the remaining time to horizon $\tau$, and the risk aversion parameter $\gamma$.
    \item \textbf{Probability Evaluation:} The theoretical Poisson intensities $\lambda^{bid}$ and $\lambda^{ask}$ are mapped into discrete transition probabilities for the interval $\Delta t$, defined as
    \begin{align*}
        P_{fill} = 1 - e^{-\lambda \Delta t}.
    \end{align*}
    \item \textbf{Execution Matching:} Two independent uniform random variables, $U^{bid}_{t_i}, U^{ask}_{t_i} \sim \mathcal{U}(0,1)$, are compared with the respective fill probabilities. If $U < P_{fill}$, the order is considered executed.
    \item \textbf{State Update:} In the event of a bid execution, inventory increases ($q_{t_{i+1}} = q_{t_i} + 1$) and the cash account decreases by the bid price ($X_{t_{i+1}} = X_{t_i} - p^{bid}$). Conversely, an execution at the ask reduces inventory and increases cash.
    \item \textbf{Mark-to-Market:} Performance is recorded by updating the total wealth (PnL), defined as
    \begin{align*}
        W_{t_{i+1}} = X_{t_{i+1}} + q_{t_{i+1}} S_{t_{i+1}}.
    \end{align*}
\end{enumerate}

\section{Experimental Control: Common Random Numbers (CRN)}
To rigorously isolate the performance differential between the symmetric benchmark and the inventory-aware policies, the experimental design employs the variance reduction technique known as Common Random Numbers (CRN).

In the Monte Carlo probability space, the entire stochastic scenario is pre-calculated via vectorized matrices. For each $m$-th simulated trajectory, a single master seed generates, at first, the sequence of Brownian increments ($\Delta W^{(m)}$) driving the mid-price process. Then the uniform random variables ($U^{bid, (m)}, U^{ask, (m)}$) governing order execution at each tick.

Both the Naive dealer and the Avellaneda-Stoikov dealer are tested on the exact same scenario. By subjecting the algorithms to the same price shocks and client order flows, the experimental design eliminates inter-path stochastic noise. Any observed disparity in terminal wealth or inventory accumulation is therefore solely attributable to the internal risk-management logic of the respective algorithms.

\section{Performance Metrics}
The simulator tracks the state variables to extract terminal distributions. The comparative analysis is based on the following key metrics:

\begin{itemize}
    \item Expected Terminal Wealth ($\mathbb{E}[W_T]$): the average profit generated across all Monte Carlo paths, representing the pure spread-capture capability.
    \item Wealth Variance ($Var(W_T)$): the dispersion of terminal PnL, acting as the primary proxy for unmanaged inventory and price risk.
    \item Daily Sharpe Ratio: calculated as the ratio between average PnL and its standard deviation, providing a standardized measure of risk-adjusted performance.
    \item Expected Absolute Terminal Inventory ($\mathbb{E}[|q_T|]$): the average magnitude of the open position remaining at the end of the trading horizon $T$. This metric quantifies the strategy's capacity for mean-reversion, measuring how effectively it avoids overnight directional exposure.
    \item Intraday Inventory Excursions ($\mathbb{E}[Std(q_t)]$): the average standard deviation of the inventory throughout the trading session. This metric captures the intraday risk profile, distinguishing between a strategy that tightly bounds its position (capital efficiency) and one that suffers massive directional swings, regardless of the final terminal state.
    \item Expected Turnover ($\mathbb{E}[\text{Turnover}]$): the average volume of executed trades per simulation, serving as a proxy for the strategy's operational efficiency and its potential vulnerability to unmodeled transaction costs.
\end{itemize}

The analysis does not rely on the calibration of real historical data but rather on the definition of stylized scenarios. This approach allows for systematic stress-testing of the models in controlled regimes (varying volatility $\sigma$ and book density $\kappa$), ensuring maximum reproducibility and isolating the structural mechanisms of the Avellaneda-Stoikov framework without the biases inherent to specific market datasets.

\chapter{Multivariate Extension: Covariance-Aware Inventory Control}
\miniabstract{The multivariate version replaces scalar inventory risk by portfolio risk. The relevant object is no longer $\sigma^2 q_t$, but the covariance pressure $\Sigma q_t$, which tells the dealer how each quote contributes to the risk of the whole book.}

PARLARE MEGLIO DELLA FATTORIZZAZIONE DI CHOLESKY PER SIMULARE IL MOTO BROWNIANO ARITMETICO MULTIVARIATO.



\section{Why Independent Single-Asset Copies Are Not Enough}
The single-asset model can be run separately for many instruments, but this is not a genuine multi-asset market maker. If prices are correlated, the risk of holding a vector of inventories is not the sum of isolated scalar risks. A long position in one asset can either amplify or hedge the risk created by another position. The practical consequence is that the dealer should not decide the bid and ask of asset $i$ using only $q_t^i$; the quote should also react to the rest of the book.

This is the main modeling step emphasized in the multi-asset treatments of market making and optimal trading. \textcite{CarteaJaimungalPenalva2015} develop the stochastic-control language used for market making, inventory penalties, and related multi-instrument trading problems. \textcite{Gueant2016} presents the Avellaneda--Stoikov model together with multi-asset generalizations, and \textcite{BergaultEvangelistaGueantVieira2022} studies closed-form approximations for multi-asset market making when direct dynamic programming becomes too expensive.

\section{State Variables and Portfolio Risk}
Let the dealer quote $d$ assets. The mid-price and inventory become vectors:
\[
S_t=(S_t^1,\ldots,S_t^d)^\top,\qquad q_t=(q_t^1,\ldots,q_t^d)^\top .
\]
For the arithmetic benchmark, the vector price process is
\[
dS_t=\mu\,dt+L\,dW_t,\qquad \Sigma=LL^\top,
\]
where $\Sigma$ is the instantaneous covariance matrix of price increments. In the data-free benchmark used here, $\mu$ is set to zero; the purpose is to isolate inventory control rather than directional alpha.

If inventory is frozen at $q_t$ over the remaining horizon $\tau=T-t$, the marked-to-market exposure is $q_t^\top(S_T-S_t)$ and its conditional variance is
\[
\mathrm{Var}_t\!\left(q_t^\top(S_T-S_t)\right)=\tau\,q_t^\top\Sigma q_t.
\]
Under CARA preferences, the corresponding certainty-equivalent inventory charge is proportional to
\[
C_t(q_t)=\frac{\gamma}{2}\tau\,q_t^\top\Sigma q_t.
\]
The scalar penalty $\frac{\gamma}{2}\tau\sigma^2 q_t^2$ is therefore replaced by a quadratic portfolio penalty. The gradient of this penalty is
\[
\nabla_q C_t(q_t)=\gamma\tau\,\Sigma q_t.
\]
This vector is the core object of the multivariate extension.

\section{Reservation Vector and Marginal Quotes}
The covariance-aware central reservation vector is
\[
m_t=S_t-\gamma\tau\,\Sigma q_t.
\]
Its $i$-th component depends on the whole inventory vector through $(\Sigma q_t)_i$. If asset $i$ is positively correlated with assets already held long, then $(\Sigma q_t)_i$ is positive and the quote center for asset $i$ moves downward: the dealer becomes less willing to buy more and more willing to sell. If asset $i$ hedges the book, the same term can be small or negative, so the dealer may quote it more aggressively even when $q_t^i$ itself is zero.

For implementation it is useful to distinguish the central reservation vector from the exact marginal bid and ask reservation prices. Let $e_i$ denote the $i$-th unit vector. The risk charge of buying one more unit of asset $i$ is
\[
C_t(q_t+e_i)-C_t(q_t)
=\gamma\tau(\Sigma q_t)_i+\frac{\gamma}{2}\tau\Sigma_{ii},
\]
while the risk charge released by selling one unit is
\[
C_t(q_t)-C_t(q_t-e_i)
=\gamma\tau(\Sigma q_t)_i-\frac{\gamma}{2}\tau\Sigma_{ii}.
\]
Therefore the marginal reservation prices used by the multi-asset code are
\[
r_{t,i}^{bid}=S_t^i-\gamma\tau(\Sigma q_t)_i-\frac{\gamma}{2}\tau\Sigma_{ii},
\qquad
r_{t,i}^{ask}=S_t^i-\gamma\tau(\Sigma q_t)_i+\frac{\gamma}{2}\tau\Sigma_{ii}.
\]
With exponential fill intensities $\lambda_i(\delta)=A_i e^{-k_i\delta}$, the liquidity half-spread for asset $i$ is
\[
\delta_i^*=\frac{1}{\gamma}\log\!\left(1+\frac{\gamma}{k_i}\right).
\]
The resulting quotes are
\[
p_{t,i}^{bid}=r_{t,i}^{bid}-\delta_i^*,
\qquad
p_{t,i}^{ask}=r_{t,i}^{ask}+\delta_i^*.
\]
This formula keeps the original Avellaneda--Stoikov separation between inventory risk and execution risk: $\Sigma q_t$ determines the skew, $\Sigma_{ii}$ contributes the marginal risk half-spread, and $k_i$ controls the liquidity half-spread.

\section{Simulation Design for the Companion Notebook}
The companion notebook implements the multivariate extension as an explicit two-asset diagnostic rather than as a full empirical study. Correlated Brownian shocks are generated by a Cholesky factor $L$ of $\Sigma$, so that each step has covariance $\Sigma\Delta t$. The same matrix-valued Gaussian shocks and the same matrix-valued uniform fill draws are reused for the naive and inventory-aware policies, preserving the common-random-number design of the single-asset simulations.

The execution model remains deliberately simple. Fills are conditionally independent across assets given the quotes, and each asset has its own $A_i$ and $k_i$. This isolates the covariance channel in the inventory term. Cross-excitation of order flow, queue priority, market impact, client tiering, and variable trade sizes are important extensions, but they would add microstructure mechanisms that are outside the current Article~1 benchmark.

The multivariate output reports terminal PnL, terminal inventory vectors, the risk norm $\sqrt{q_T^\top\Sigma q_T}$, and the covariance pressure $\Sigma q_T$. These diagnostics answer a different question from the scalar Monte Carlo table: not only whether inventory is smaller asset by asset, but whether the remaining book is less risky after accounting for correlations.

\chapter{Implementation and Validation}
\miniabstract{The simulation can be implemented in Python with a modular design separating price dynamics, execution logic, quoting policy, and risk controls. Validation relies on sanity checks that should hold regardless of the chosen diffusion case.}

\section{General Simulation Schema}
\begin{algorithm}[H]
\caption{Event-driven market making simulation}
\DontPrintSemicolon
\KwIn{Model parameters, horizon, quoting policy, and random seed}
\KwOut{Paths for $S_t$, $q_t$, $X_t$, $W_t$, and summary metrics}
Initialize $S_0$, $q_0$, and $X_0$\;
\For{each event time or time step until the horizon ends}{
  update the diffusion state\;
  compute bid and ask quotes from the policy\;
  evaluate fill intensities\;
  simulate executions\;
  update inventory and cash\;
  record performance metrics and diagnostics\;
}
\Return{Monte Carlo statistics and path summaries}\;
\end{algorithm}

\section{Placeholder Quoting Function}
\begin{lstlisting}[caption={Illustrative structure of an Avellaneda--Stoikov quoting rule},label={lst:as_policy_1}]
import numpy as np

def intensita(delta, A, k):
    return A * np.exp(-k * delta)

def quoting_policy_as(S, q, t, T, sigma, gamma, A, k):
    tau = T - t
    r = S - q * gamma * sigma**2 * tau
    spread = 2.0 / k + gamma * sigma**2 * tau
    delta_ask = spread / 2.0 + (r - S)
    delta_bid = spread / 2.0 - (r - S)
    p_bid = S - max(delta_bid, 0.0)
    p_ask = S + max(delta_ask, 0.0)
    return p_bid, p_ask
\end{lstlisting}

\section{Sanity Checks}
Expected sanity checks include symmetry at $q=0$, monotonic widening of the spread as $\gamma$ or $\sigma$ increases, stronger inventory pressure over the trading session in finite-horizon cases, and covariance-consistent quote shifts in the multivariate diagnostic. These checks are especially useful because the analysis does not rely on market data and must therefore establish internal consistency by design.

\chapter{Conclusion}
\miniabstract{The paper provides the classical baseline against which more elaborate market making models can later be compared. Its main value lies in keeping the execution engine fixed while varying only the mid-price model and then extending the same inventory logic to a covariance-aware multivariate setting.}

The Avellaneda--Stoikov framework is particularly useful because it formalizes in a compact way the central trade-off of market making: spread capture versus inventory and execution risk. Within a single-asset setting, arithmetic Brownian motion and geometric Brownian motion already generate meaningfully different quoting environments. The multivariate extension then shows how the same inventory-aware logic changes when the dealer manages a vector of correlated positions. The comparison remains coherent because all cases share the same accounting, execution mechanics, and inventory-aware control principle.

The main expectation is that inventory-based quoting will reduce inventory excursions and stabilize wealth relative to naive symmetric quoting, though the exact strength of the effect will vary with volatility, risk aversion, and covariance structure. This establishes a clean numerical benchmark for later extensions involving richer volatility dynamics, more realistic order flow, and production-grade multi-asset inventory control.

\appendix
\chapter{Appendix: Additional Technical Notes}
The HJB derivations, numerical approximations, and alternative state-dependent specifications can be developed further in a later version of the paper. The present draft keeps the analytical exposition intentionally focused on the benchmark mechanisms required for simulation and interpretation.

\begin{summary}
\textbf{Objective:} Build a finite-horizon single-asset benchmark and a covariance-aware multi-asset extension for high-frequency market making under the Avellaneda--Stoikov framework.\\
\textbf{Method:} Event-driven Monte Carlo simulation with a common execution engine; mid-price model varied between ABM and GBM; multivariate diagnostic driven by a covariance matrix.\\
\textbf{Findings:} Inventory-aware quoting should reduce inventory excursions, lower the dispersion of marked-to-market wealth, and clarify how covariance-aware inventory pressure changes quote centers in the multivariate setting.\\
\textbf{Risks:} The baseline abstracts from queue position, adverse selection, and explicit latency effects.\\
\textbf{Recommendation:} Use this paper as the clean benchmark before moving to rough volatility, Hawkes flow, and production-grade multivariate inventory control.
\end{summary}

\begin{references}
\bibitem{AvellanedaStoikov2008}
Marco Avellaneda and Sasha Stoikov,
\emph{High-frequency trading in a limit order book},
\emph{Quantitative Finance} 8(3), 217--224, 2008.

\bibitem{BergaultEvangelistaGueantVieira2022}
Philippe Bergault, David Evangelista, Olivier Gu\'eant, and Douglas Vieira,
\emph{Closed-form approximations in multi-asset market making},
arXiv:1810.04383v5, 2022.

\bibitem{CarteaJaimungalPenalva2015}
\'Alvaro Cartea, Sebastian Jaimungal, and Jos\'e Penalva,
\emph{Algorithmic and High-Frequency Trading},
Cambridge University Press, 2015.

\bibitem{GouldEtAl2013}
Martin D. Gould, Mason A. Porter, Stuart Williams, Mark McDonald, Daniel J. Fenn, and Sam D. Howison,
\emph{Limit order books},
\emph{Quantitative Finance} 13(11), 1709--1742, 2013.

\bibitem{GueantLehalleFernandezTapia2013}
Olivier Gu\'eant, Charles-Albert Lehalle, and Joaquin Fernandez-Tapia,
\emph{Dealing with the inventory risk: a solution to the market making problem},
\emph{Mathematics and Financial Economics} 7(4), 477--507, 2013.

\bibitem{Gueant2016}
Olivier Gu\'eant,
\emph{The Financial Mathematics of Market Liquidity: From Optimal Execution to Market Making},
CRC Press, 2016.

\bibitem{Harris2003}
Larry Harris,
\emph{Trading and Exchanges: Market Microstructure for Practitioners},
Oxford University Press, 2003.

\bibitem{Hasbrouck2007}
Joel Hasbrouck,
\emph{Empirical Market Microstructure: The Institutions, Economics, and Econometrics of Securities Trading},
Oxford University Press, 2007.

\bibitem{OHara1995}
Maureen O'Hara,
\emph{Market Microstructure Theory},
Blackwell, 1995.
\end{references}

\end{document}
